\documentclass[11pt,letterpaper]{article}
\usepackage[margin=1in]{geometry}
\usepackage{amsmath,amssymb,booktabs,tabularx,hyperref,xurl}
\hypersetup{hidelinks,pdfauthor={Shokoofeh Akbari, Oleg Shupliakov, Yousef Jamali}}
\begin{document}
\begin{center}
{\Large\bfseries S1 Text. Supporting methods and numerical results}\par
\medskip
Electrostatic--hydrophobic cooperation in the mesoscale assembly of synapsin and synaptic vesicles\par
\medskip
Shokoofeh Akbari, Oleg Shupliakov, and Yousef Jamali
\end{center}
\section*{S1. Design and structural observables}
The factorial design contains four E/H conditions and three seed blocks,
9310002--9310004. Each trajectory has 200 SVs and 1,807 Syn agents in a
spherical compartment of volume $0.1\,\mu\mathrm m^3$. All trajectories
were followed to 84 ms at a fixed integration step of 0.0007 ms.
Van der Waals interactions remain enabled in every condition.
The 600 unique saved frames are separated by 0.14 ms. Continuations
retain the same physical parameters and saved random state; identical
checkpoint frames are counted once.

The principal 84-ms summaries average frames 421--600 (58.94--84 ms)
within each trajectory, followed by a mean and sample SD across seeds.
The earlier 28-ms summaries use frames 141--200 (19.74--28 ms).
Both apply the same final-30\% rule. Each trajectory is one independent
replicate. The two horizons are repeated observations of the same runs.

The association graph uses $\|x_i-x_j\|\leq a_i+a_j+2\,\mathrm{nm}$,
where $a_i$ is the outer radius. $L$ is the all-agent largest-component
fraction, $Q_{\rm Syn}$ is the non-isolated Syn fraction and $L_{\rm SV}$
is the fraction of all SVs in the component selected by total agent count.
Non-isolated Syn includes Syn--Syn groups and Syn--SV associations.
The original machine-readable column ``synapsin\_recruited\_fraction''
and the updated column ``syn\_nonisolated\_fraction'' both mean $Q_{\rm Syn}$.

\section*{S2. Complete per-run structural observables at 84 ms}
\begin{center}\small
\begin{tabular}{llrrrr}
\toprule
Condition & Seed & $L$ & $Q_{\rm Syn}$ & $L_{\rm SV}$ & Drift-limited \\ \midrule
E0H0 & 9310002 & 0.027675 & 0.318434 & 0.084889 & 0.000001 \\
E0H0 & 9310003 & 0.054091 & 0.319935 & 0.171139 & 0.000000 \\
E0H0 & 9310004 & 0.028647 & 0.313272 & 0.086889 & 0.000000 \\
E0H1 & 9310002 & 0.337812 & 0.560004 & 0.545417 & 0.001161 \\
E0H1 & 9310003 & 0.413553 & 0.592714 & 0.587556 & 0.002336 \\
E0H1 & 9310004 & 0.280632 & 0.550593 & 0.435806 & 0.001121 \\
E1H0 & 9310002 & 0.445244 & 0.759017 & 0.484389 & 0.000000 \\
E1H0 & 9310003 & 0.599812 & 0.779850 & 0.616083 & 0.000000 \\
E1H0 & 9310004 & 0.462299 & 0.752702 & 0.496528 & 0.000000 \\
E1H1 & 9310002 & 0.921386 & 0.953849 & 0.908056 & 0.005209 \\
E1H1 & 9310003 & 0.912982 & 0.949324 & 0.868333 & 0.005391 \\
E1H1 & 9310004 & 0.925491 & 0.951593 & 0.896361 & 0.005518 \\
\bottomrule\end{tabular}\end{center}
The final column gives the whole-84-ms mean fraction of particle updates
affected by the deterministic drift limiter, weighted by integration
steps across segments. Brownian increments are retained in full.
All segments completed with finite positions and no unconverged
projection steps at the original 0.005-nm residual tolerance.
\clearpage
\section*{S3. Paired interaction contrasts at both horizons}
The contrast is $I_s=Y_{11,s}-Y_{10,s}-Y_{01,s}+Y_{00,s}$.
The following table retains every seed and each structural observable.
\begin{center}\small
\begin{tabular}{llrrr}
\toprule
Horizon (ms) & Seed & $I_L$ & $I_{Q_{\rm Syn}}$ & $I_{L_{\rm SV}}$ \\ \midrule
28 & 9310002 & 0.359043 & -0.029266 & 0.380500 \\
28 & 9310003 & 0.313237 & -0.031452 & 0.307500 \\
28 & 9310004 & 0.444503 & -0.026075 & 0.455250 \\
84 & 9310002 & 0.166005 & -0.046738 & -0.036861 \\
84 & 9310003 & -0.046291 & -0.103305 & -0.164167 \\
84 & 9310004 & 0.211208 & -0.038431 & 0.050917 \\
\bottomrule\end{tabular}\end{center}
At 84 ms, the mean $I_L$ is 0.110307, with between-seed SD 0.137488. The approximate 95\% paired Student-$t$ interval is [-0.231233, 0.451848].

The main text uses paired differences and sample SDs. There are three
seed blocks and two degrees of freedom for paired intervals. Normality
of seed differences is difficult to assess with three blocks. The
full-precision conditional contrasts are retained in S1 Data. The
original 28-ms statistical audit remains a historical analysis of that
horizon and is stored separately from the 84-ms tables.

\section*{S4. Association-shell sensitivity at 84 ms}
Only the graph threshold changes; trajectories and forces are fixed.
Entries are means of three per-run late-window means.
\begin{center}\small
\begin{tabular}{lrrr}
\toprule
Condition & Shell (nm) & $L$ & $L_{\rm SV}$ \\ \midrule
E0H0 & 0 & 0.018471 & 0.070759 \\
E0H0 & 2 & 0.036805 & 0.114306 \\
E0H0 & 4 & 0.073231 & 0.193611 \\
E0H0 & 6 & 0.143276 & 0.328630 \\
E0H1 & 0 & 0.257311 & 0.402796 \\
E0H1 & 2 & 0.343999 & 0.522926 \\
E0H1 & 4 & 0.411524 & 0.614833 \\
E0H1 & 6 & 0.473130 & 0.692194 \\
E1H0 & 0 & 0.346878 & 0.363213 \\
E1H0 & 2 & 0.502452 & 0.532333 \\
E1H0 & 4 & 0.596171 & 0.644713 \\
E1H0 & 6 & 0.666221 & 0.730852 \\
E1H1 & 0 & 0.887193 & 0.857648 \\
E1H1 & 2 & 0.919953 & 0.890917 \\
E1H1 & 4 & 0.929835 & 0.909870 \\
E1H1 & 6 & 0.938186 & 0.927667 \\
\bottomrule\end{tabular}\end{center}
\clearpage
\section*{S5. Species-resolved contact analysis at 84 ms}
For each Syn agent $i$, the SV-neighbor count is
\[
d_i^{\rm SV}=\sum_{j\in\mathcal V}\mathbf1\{\|x_i-x_j\|\leq a_i+a_j+2\,\mathrm{nm}\}.
\]
The two Syn fractions count $d_i^{\rm SV}\geq1$ and $d_i^{\rm SV}\geq2$,
respectively, divided by $N_{\rm Syn}$. Reciprocal SV coverage counts
SVs with at least one Syn neighbor, divided by $N_{\rm SV}$.
The same threshold is applied to every condition. Values are averages
over 58.94--84 ms within each run.
\begin{center}\small
\begin{tabular}{llrrr}
\toprule
Condition & Seed & Syn: $\geq1$ SV & Syn: $\geq2$ SV & SV: $\geq1$ Syn \\ \midrule
E0H0 & 9310002 & 0.191407 & 0.012384 & 0.805500 \\
E0H0 & 9310003 & 0.191063 & 0.013552 & 0.810472 \\
E0H0 & 9310004 & 0.189341 & 0.012141 & 0.814444 \\
E0H1 & 9310002 & 0.454052 & 0.124909 & 0.957778 \\
E0H1 & 9310003 & 0.501024 & 0.192194 & 0.950778 \\
E0H1 & 9310004 & 0.438929 & 0.119185 & 0.952000 \\
E1H0 & 9310002 & 0.745207 & 0.193448 & 0.993083 \\
E1H0 & 9310003 & 0.768591 & 0.253127 & 0.988639 \\
E1H0 & 9310004 & 0.738695 & 0.182454 & 0.992444 \\
E1H1 & 9310002 & 0.951928 & 0.664957 & 0.984139 \\
E1H1 & 9310003 & 0.947703 & 0.676736 & 0.975222 \\
E1H1 & 9310004 & 0.949920 & 0.664690 & 0.983000 \\
\bottomrule\end{tabular}\end{center}
Mean unordered edge counts per snapshot:
\begin{center}\small
\begin{tabular}{llrrr}
\toprule
Condition & Seed & Syn--Syn & Syn--SV & SV--SV \\ \midrule
E0H0 & 9310002 & 157.572 & 368.650 & 105.744 \\
E0H0 & 9310003 & 159.033 & 370.122 & 122.639 \\
E0H0 & 9310004 & 151.844 & 364.400 & 102.083 \\
E0H1 & 9310002 & 398.800 & 1069.344 & 179.533 \\
E0H1 & 9310003 & 610.489 & 1321.411 & 208.717 \\
E0H1 & 9310004 & 398.356 & 1031.139 & 164.022 \\
E1H0 & 9310002 & 296.000 & 1696.178 & 28.694 \\
E1H0 & 9310003 & 366.533 & 1846.317 & 39.150 \\
E1H0 & 9310004 & 281.094 & 1664.517 & 27.378 \\
E1H1 & 9310002 & 1331.250 & 2982.600 & 126.700 \\
E1H1 & 9310003 & 1390.100 & 3020.233 & 132.522 \\
E1H1 & 9310004 & 1396.917 & 2953.311 & 125.256 \\
\bottomrule\end{tabular}\end{center}
\clearpage
\section*{S6. Short-lag mobility and experimental scale reference}
Syn positions are projected onto $xy$, converted from nm to $\mu$m,
and fitted using per-agent time-averaged MSD with a free intercept over
native lags 2.10--3.92 ms (14 lag points). The apparent coefficient is
slope/4 in $\mu\mathrm m^2/\mathrm s$. Negative fitted slopes are retained.
The full 600-frame trajectory is used for the main figure. The final
180 frames provide a window-sensitivity comparison.
\begin{center}\small
\begin{tabular}{llrrrr}
\toprule
Condition & Seed & $D$: full & $D$: late & Negative: full & Negative: late \\ \midrule
E0H0 & 9310002 & 0.141098 & 0.143295 & 0.0000 & 0.0033 \\
E0H0 & 9310003 & 0.140255 & 0.137813 & 0.0000 & 0.0033 \\
E0H0 & 9310004 & 0.139690 & 0.139443 & 0.0000 & 0.0050 \\
E0H1 & 9310002 & 0.131435 & 0.129700 & 0.0000 & 0.0039 \\
E0H1 & 9310003 & 0.129508 & 0.122338 & 0.0000 & 0.0017 \\
E0H1 & 9310004 & 0.131327 & 0.128424 & 0.0000 & 0.0022 \\
E1H0 & 9310002 & 0.110672 & 0.108970 & 0.0000 & 0.0033 \\
E1H0 & 9310003 & 0.110623 & 0.104226 & 0.0000 & 0.0061 \\
E1H0 & 9310004 & 0.110988 & 0.107553 & 0.0000 & 0.0061 \\
E1H1 & 9310002 & 0.062747 & 0.055275 & 0.0000 & 0.0100 \\
E1H1 & 9310003 & 0.063020 & 0.053947 & 0.0000 & 0.0083 \\
E1H1 & 9310004 & 0.062057 & 0.055238 & 0.0006 & 0.0188 \\
\bottomrule\end{tabular}\end{center}
Hoffmann et al. (Nature Communications 2023, DOI
\href{https://doi.org/10.1038/s41467-023-42372-6}{10.1038/s41467-023-42372-6})
provide a Fig. 2h short-time mobility scale: source-data means are
0.0189192370 with SV and 0.0253090377 without SV, in
$\mu\mathrm m^2/\mathrm s$. Source data:
\url{https://doi.org/10.6084/m9.figshare.24179445.v1}.
Their experiment selects condensate-core tracks at a 1-ms cadence,
whereas the model analysis includes all Syn at a 0.14-ms saved cadence.
Simulation positions have no localization error or exposure blur, and
every model condition contains SV. Effective friction was chosen with
reference to the literature before this comparison. The comparison provides a mobility scale reference, with higher absolute
mobility in the model. Experimental uncertainty was not recalculated.


\section*{S7. Reproduction and software checks}
S1 Data contains full-precision inputs for all main figures, per-run
structural and mobility values, shell comparisons and paired contrasts.
The original 28-ms tables remain separate. The updated repository adds
12 common-84-ms analysis-input files, preserving all 600 frames and the
exact stored values and types of the included arrays. Original raw
outputs remain preserved separately. The portable script
\nolinkurl{analysis/common84ms.py} reproduces the new tables, and
\nolinkurl{scripts/render_common84ms.py} generates the four figures.

Raw-source hashes and exact checkpoint joins were checked, including
positions, identities, radii and random-number state. Reconstruction of
the earlier 2,400 frames reproduced the original structural and contact
tables to absolute tolerance $10^{-12}$. Selected graph calculations
were checked against all-pairs enumeration. Saved positions and radii
use float32; one reconstructed overlap residual differed from the
solver tolerance by $1.52\times10^{-6}$ nm, within its rounding bound.
The recorded double-precision solver residual remained below 0.005 nm.
Detailed numerical checks are retained with the analysis data.

\section*{S8. Later assembly trends}
For each trajectory, a straight line was fitted to $L$ within each
14-ms period. Entries are means of three per-run means and slopes.
The periods below show the later development of each condition.
\begin{center}\small
\begin{tabular}{llrr}
\toprule
Condition & Period (ms) & Mean $L$ & Slope (ms$^{-1}$) \\ \midrule
E0H0 & 42--56 & 0.031581 & 0.000520 \\
E0H0 & 56--70 & 0.040241 & -0.000322 \\
E0H0 & 70--84 & 0.034406 & -0.000352 \\
E0H1 & 42--56 & 0.191023 & 0.003080 \\
E0H1 & 56--70 & 0.297030 & 0.008990 \\
E0H1 & 70--84 & 0.370037 & 0.005398 \\
E1H0 & 42--56 & 0.318683 & 0.007319 \\
E1H0 & 56--70 & 0.423941 & 0.006971 \\
E1H0 & 70--84 & 0.560191 & 0.006058 \\
E1H1 & 42--56 & 0.831795 & 0.008818 \\
E1H1 & 56--70 & 0.910580 & 0.001990 \\
E1H1 & 70--84 & 0.925225 & 0.000381 \\
\bottomrule\end{tabular}\end{center}
\end{document}
